\section*{Cycle Fan Predicate: syntax and semantic boundary}

\noindent\textbf{Status: OPEN for arbitrary fixed $k,R,m$.}
The previously displayed expression containing
$\exists^{\ge m}z$ was described as an ordinary $\mathrm{FO}^k$
formula without proving that it can be expressed using at most $k$
variable names. As written, $\exists^{\ge m}$ is counting-quantifier
notation; counting logic $\mathrm{C}^k$ must not be silently identified
with ordinary $\mathrm{FO}^k$.

For contrast, with at least $m+1$ available variable names the
\emph{simple} property that $x$ has $m$ distinct neighbors is defined in
ordinary first-order logic by
\[
\psi_m(x):=\exists z_1\cdots\exists z_m\,
\Bigl(\bigwedge_{i=1}^{m}E(x,z_i)\ \land\
\bigwedge_{1\le i<j\le m}z_i\ne z_j\Bigr).
\]
This uses $m+1$ names including $x$, and says nothing about independent
cycles or a shared fan. It is \emph{not} a replacement proof for the
cycle-fan claim. Adding bounded paths, fixed anchors, and cycle
independence requires a separate formal construction with its variable
budget and semantic assumptions explicit.

\noindent\textbf{OPEN obligations:} Fix the graph vocabulary and logic;
define an $(m,R)$-fan and independent cycles; construct a corresponding
formula with an audited variable bound; then prove any proposed
cycle-rank-to-fan implication independently. None of those steps follows
from the existing counting-quantifier display.
